\chapter{Conclusions and Future Work}
\label{ch:conclusions}

\section{Summary of Contributions}

This project set out to test a narrow architectural claim: that adding one
backward edge to a retrieval-augmented pipeline changes what the system is
capable of, and that the change can be measured rather than merely described.

What was built is a four-agent system coordinated by an explicit state machine.
The Planner decomposes a question into a validated directed acyclic graph of
sub-queries connected by placeholders. The Retrieval Agent routes each sub-query
to BM25, dense, or hybrid search via a seven-rule decision table over lexical
features --- of which, on the evaluated data, only the first rule ever fired,
because the Planner attached a tool hint to every sub-query --- fuses multiple
formulations with reciprocal rank
fusion~\cite{cormack2009reciprocal}, and reranks only when the fused margin is
narrow. The Knowledge Graph Navigator links entities by alias matching and runs
bidirectional breadth-first search over an entity graph induced from corpus
passage text alone. The Verifier resolves citations, scores entailment with a
DeBERTa cross-encoder, and --- below its confidence threshold --- returns a
\texttt{ReplanDirective} to the Planner instead of an answer. The four agents
cover all five tasks of Table~1 of the brief, and the system runs on one
consumer GPU with no paid API.

Three properties are established independently of the benchmark outcome,
because they were measured during construction. Structured planning with a local
8B model is reliable when the output is schema constrained: across the planning
calls measured in development the model produced correctly shaped decompositions
in 8 of 8 cases with zero format retries. The same model is unreliable at
describing what it produced, mislabelling its own plan's \texttt{strategy} field
in 7 of those 8, which is why strategy is derived from the DAG shape and the
model's label discarded. And agency is affordable here only because most of it
is deterministic: a full agentic question costs at worst 15 model calls against
a cap of 20 and typically four to seven, since tool routing, entity linking,
graph traversal, and the confidence arithmetic consume none.

\section{Empirical Findings}

Every number in this section is taken from the tables of
Chapter~\ref{ch:implementation}, which are generated from the run artefacts;
none is typed by hand, and where a run has not finished the finding is withheld
rather than estimated. The headline is stated once, plainly, before the
structure: the agentic system is not the most accurate system in this study, it
is the only one that can show its work, and those two facts are the result.

The evaluation answers four questions. Each is stated explicitly because each
has a defensible negative answer, and the protocol was fixed before any of them
was resolved.

\textbf{Does decomposition beat single-shot retrieval?} Two reference points
are needed, because no single baseline is comparable on every column.
\texttt{hybrid\_rerank}, the strongest non-agentic ranker, is the reference
for retrieval and supporting-fact quality; it attempts no answer, so it cannot
be the reference for answer quality, and the generated tables render its
answer columns as \texttt{--} rather than as a zero that would read as a
failed attempt. The answer reference is \texttt{self\_ask}, the strongest
baseline that produces an answer string, and every answer-quality delta below
is a paired bootstrap against it over the identical 250 questions. An earlier
draft of this report described the planner's fallback floor as
\texttt{hybrid\_rerank} ``by construction'' and argued that the agentic system
could lose only through judgement; the ablation grid shows the floor is a
one-node agentic system that still synthesises and verifies, and the argument
is withdrawn.
The answer is no, and it is a significant no on one of the two datasets. On
HotpotQA \texttt{agentic\_full} scores EM 0.432 and F1 0.510 where
\texttt{self\_ask} scores 0.504 and 0.603 and \texttt{naive\_rag} 0.496 and
0.581; the paired difference in F1 against \texttt{self\_ask} is $-0.093$ with
a 95\% interval of $[-0.145, -0.039]$, which excludes zero, so the deficit is
established and not a fluctuation. On 2WikiMultihopQA the figures are 0.224
and 0.267 against \texttt{self\_ask}'s 0.264 and 0.310 and
\texttt{naive\_rag}'s 0.152 and 0.200, and the paired difference of $-0.043$
$[-0.100, +0.014]$ spans zero, so there the deficit is a consistent point
estimate and nothing stronger. Against single-shot retrieve-then-generate the
picture is more favourable on the harder benchmark, where 2WikiMultihopQA's
longer chains punish a single query and \texttt{agentic\_full} sits seven
exact-match points above \texttt{naive\_rag}, but no paired interval between
those two is computed and their deltas against the common reference overlap,
so that ordering is not certified either. What the grid does establish is
that decomposition with verification loses to decomposition without it on the
answer string, significantly on HotpotQA and directionally on
2WikiMultihopQA.

On evidence the answer is different in kind rather than simply reversed.
Supporting-fact F1 is 0.481 against 0.412 for \texttt{hybrid\_rerank} on
HotpotQA and 0.421 against 0.327 on 2WikiMultihopQA, both significant under
the paired bootstrap. The generative baselines were reported in an earlier
draft at 0.000 on that metric, as systems that ``emit no supporting facts'';
that was a harness defect in resolving their citations, and corrected they
score 0.405 and 0.412 on HotpotQA and 0.318 and 0.327 on 2WikiMultihopQA ---
starred in the tables, because they are scored on their whole evidence pool
where the agentic system is scored on the sentences it cited. The two
protocols trade precision against recall in opposite directions --- on
HotpotQA the agentic cited set scores a precision of 0.654 and a recall of
0.406 where \texttt{self\_ask}'s pool scores 0.302 and 0.696 --- and one F1
across them does not rank the systems. Scored the same way for every row,
on the recall of the whole evidence pool, the agentic system leads every
baseline significantly, 0.793 against 0.706 for \texttt{hybrid\_rerank} on
HotpotQA and 0.684 against 0.622 on 2WikiMultihopQA. What is categorical is
citation grounding: 0.960 and 0.976 for the agentic system against 0.000 for
every baseline, whose citations resolve to nothing. Fused
retrieval quality moves the other way --- nDCG@10 of 0.763 and 0.723 against
0.835 and 0.756 for the hybrid baseline, both significantly worse --- because
the reported ranking for a multi-hop system is the fusion of several narrow
sub-query lists and fusion dilutes the head of the list. On 2WikiMultihopQA
the two systems tie exactly on Recall@10 at 0.727, which locates the loss in
ordering rather than in coverage. So decomposition improves which evidence is
cited, and how checkably, and degrades where in a single fused list the
evidence appears; the aggregate ranking metric alone would have concealed
that.

\textbf{Does the backward edge pay for itself?} This is the central claim, and
it is measured three ways. First, the marginal effect: \texttt{agentic\_full}
against \texttt{agentic\_no\_verifier}, which removes the loop and holds
everything else constant. Second, conditional effectiveness, restricted to the
questions where a re-plan fired and comparing the first cycle's answer against
the one finally selected --- because if re-planning is rare the marginal effect
will be small even when the mechanism works perfectly on the cases it reaches.
Third, the counterweight: false rejects, where the Verifier sent back an answer
that was already correct. A loop that fires often and improves nothing, or one
that rarely fires, is a legitimate and reportable outcome; the trace records the
selected cycle per question precisely so this can be settled from evidence.
All three measurements are now available on HotpotQA, and the decisive one
answers the question the report was written to ask. The
loop is genuinely exercised: the re-plan rate is 0.472 on HotpotQA and 0.392 on
2WikiMultihopQA, so on roughly two questions in five the Verifier declined to
answer and sent a directive back to the Planner. Whatever else is true, the
backward edge is not decorative, and the anti-repetition ban list on sub-query
text is doing real work at that frequency. The counterweight is small: the
traces record 2 false rejects on HotpotQA and 5 on 2WikiMultihopQA, and those
are exactly the questions marked recoverable, meaning some discarded cycle had
been right. Selecting the answer by maximum confidence over all cycles, with
ties broken toward the earlier one, is therefore doing what it was designed to
do --- a bad re-plan almost never overwrites a good first attempt.

The failure lies in the other direction, and it is large. The Verifier accepted
92 answers on HotpotQA and 154 on 2WikiMultihopQA that the gold supporting
facts do not support. A mechanism that fires on 47.2\% of questions while
rubber-stamping that many unsupported answers is not too aggressive but far too
permissive, and the calibration evidence below says why raising the threshold
is not an easy remedy.

The marginal effect is the central result of this report, and on HotpotQA it
is positive and significant. \texttt{agentic\_no\_verifier} disables the
backward edge and holds everything else fixed; over the same 250 questions it
scores an exact match of 0.396 and F1 of 0.474 against the full system's 0.432
and 0.510, a paired $\Delta$F1 of $-0.037$ with a 95\% interval of
$[-0.065, -0.011]$ and $p = 0.008$. The loop that fires on nearly half the
questions does improve the aggregate, by just under four F1 points, at a cost
of roughly doubling the model calls, 4.26 against 2.06, and the median
latency, 28.3\,s against 16.5\,s. That is the claim the architecture was built
to test, it is measured rather than asserted, and it lands in the direction
the design predicted. It also has to be held at its true size. It is a
single-dataset result, and the report's own standard is that one dataset does
not settle a direction; the 2WikiMultihopQA ablation has not been run and is
the measurement that would make the finding a general one. And it is a
statement about the edge, not about the system: a verifier that adds four
points to an answer accuracy that is itself nine points behind
\texttt{self\_ask} has proved that feedback helps, not that this system is
the one to deploy.

The planner ablation, complete on HotpotQA at the same time, produced the
result the grid was not designed to expect. \texttt{agentic\_no\_planner}
scores \emph{above} the full system, at an exact match of 0.456 and F1 of
0.543, with $\Delta$F1 $= +0.033$ $[-0.017, +0.086]$ and $p = 0.186$ --- not
significant, but pointing against the component the brief lists first --- at
1.35 model calls and a median of 9.7\,s, roughly a third of the full system's
cost. It is not the \texttt{hybrid\_rerank} floor: it still synthesised and
verified, grounded its citations at 0.936, and its Verifier triggered a
re-plan on 18.8\% of questions, though every such re-plan was discarded as a
duplicate of the identity plan and none executed. The reading that
fits every other row is that HotpotQA's two-hop questions are largely within
reach of one well-formed hybrid query, so decomposition dilutes a ranking that
was already good and asks an 8B model to compose across sub-answers it would
not otherwise have had to. Whether the Planner earns its cost on the longer
chains of 2WikiMultihopQA is the open question that dataset's ablation row
will answer.

The graph ablation completes the HotpotQA grid with a third result that
qualifies the architecture rather than confirming it. \texttt{agentic\_no\_kg}
scores an exact match of 0.404 and F1 of 0.501, a paired $\Delta$F1 of
$-0.009$ $[-0.048, +0.029]$ with $p = 0.646$, and its retrieval is
statistically indistinguishable from the full system's --- Recall@10 of 0.820
against 0.812, nDCG@10 of 0.766 $[0.737, 0.793]$ against 0.763 $[0.734,
0.791]$ --- with fifteen bridge-link failures rather than seventeen. The
knowledge graph was added to find bridge entities, and on this dataset it
cannot be shown to find any the retriever had not already surfaced, while
accounting for almost half of the system's tool calls and more than half of
its median latency. The finding is not that graph navigation is useless; it is
that two-hop questions do not exercise it, and that the case for the component
has to be made on 2WikiMultihopQA or not at all.

\textbf{What does agency cost?} Quality is reported alongside model calls per
question, calls avoided by deterministic rules, tool invocations, plan depth,
re-plan rate, and latency, so no accuracy gain is quoted without its compute
price. On the same axis the tool-routing table is compared against model-driven
selection on a held-out slice, by agreement rate, retrieval quality, and wall
clock: the claim that a rule matches a model's decision at zero cost is only
worth making if the agreement rate is stated.
It costs about three times the model calls of the strongest baseline and rather
more than that in wall clock. \texttt{agentic\_full} spends 4.26 model calls
per question on HotpotQA and 3.80 on 2WikiMultihopQA against 1.24 and 1.46 for
\texttt{self\_ask}, and issues 14.83 and 16.00 tool calls against 2.41 and
2.86. Latency is reported as a median because one HotpotQA question ran for
18.3 hours inside a single call, an outlier that an earlier draft's mean
reported as the per-question cost. On the medians, \texttt{agentic\_full}
takes 28.3\,s per question on HotpotQA against 4.5\,s for \texttt{self\_ask}
and 4.9\,s for \texttt{hybrid\_rerank}, and 35.4\,s on 2WikiMultihopQA against
13.2\,s and 8.3\,s. Set against a result that is \emph{worse} on answer
accuracy, this is not a good trade for a system whose deliverable is an answer
string, and the report says so rather than presenting the attribution gain as
though it settled the arithmetic. It is a defensible trade only where an
unattributable answer has little value.

The \textit{Saved} column is the cost finding that had to be revised most.
An earlier draft reported 7.10 and 7.19 model calls per question displaced by
deterministic rules, more than the system spent, and called it the clearest
engineering result in the work. The count was inflated: it credited a
knowledge-graph entity-linking call on every sub-query that this
configuration never issues, and it missed the extraction-ladder rungs that
did avoid a call. Under the corrected definition, in which a saved call is
one a configured model path would have made and a rule pre-empted, the
column reads 2.60 and 3.10 calls per question against 4.26 and 3.80 spent,
and Chapter~\ref{ch:implementation} prints the per-rule decomposition beside
it. The deterministic components displace a real
minority of the inference a fully model-driven version of this design would
spend; they do not displace more than the three reasoning agents consume, and
that sentence is withdrawn wherever the earlier draft made it.

The comparison that would have tested the routing rules directly has not been
run, and as the code stands cannot be. The rule-versus-model experiment
described in Chapter~\ref{ch:methodology} --- clearing the
\texttt{heuristic\_shortcut} flag and re-running the calibration slice through
the selection prompt --- produced no artefact in \texttt{results/}, so no
agreement rate is available and none is reported. The trace also shows why
the experiment would not have measured what it was designed to measure: the
Planner attached a tool hint to every sub-query in the headline run, the
retriever honours a hint before it consults the shortcut flag, and so the
model router is unreachable and the six lexical rules behind the hint rule
never executed. The routing decision was made inside the planning call, and
what the table saved was a second call for the same decision. That is a
finding about where the division of labour between Planner and Retrieval
Agent actually fell, and it should be read as one; the claim that lexical
rules decide as a model would have decided is untested.

\textbf{Where does it fail?} The error analysis partitions failures by the stage
that caused them --- decomposition, first-hop retrieval, placeholder resolution,
extraction, synthesis --- because an aggregate score attributes nothing, and a
system that fails mostly at retrieval calls for entirely different subsequent
work than one that fails mostly at composition.
It fails at composition, and it used to fail at retrieval. Measured against
\texttt{hybrid\_rerank}, the full system reduces the share of questions
labelled \texttt{decomposition\_error} from 0.248 to 0.064 on HotpotQA and from
0.552 to 0.088 on 2WikiMultihopQA, and \texttt{bridge\_link\_failure} from
0.552 to 0.068 and from 0.200 to 0.028. Those two categories are precisely the
ones Chapter~\ref{ch:introduction} argued a single-shot ranker cannot address
in principle, and the machinery built to address them has very largely done so.
Neither dataset records a single parse failure or budget exhaustion in any
configuration, so nothing in these results is an artefact of malformed model
output or of a question that ran out of allowance, and only 16 of 142 HotpotQA
errors and 23 of 194 on 2WikiMultihopQA fall through the taxonomy unlabelled.

What is left is \texttt{synthesis\_error}, at 0.228 and 0.236 --- the dominant
agentic failure mode on both datasets and more than three times the
corresponding rate for \texttt{self\_ask}, 0.064 and 0.076. Read together with
a supporting-fact F1 that leads every retrieval baseline and an exact match
that does not, the finding is specific: the system retrieves the right
sentences and then extracts or composes the wrong span from them. This is the
single most useful thing the evaluation produced, because it says where the
next unit of engineering effort belongs. It does not belong in more planning
--- the planner ablation above says as much --- nor in deeper traversal or
better fusion, all of which are already delivering the evidence. It belongs
in answer extraction, and by extension in a verifier that can distinguish a
well-cited answer from a correct one --- a distinction the present confidence
blend, which scores an answer \emph{sentence} against its citations and never
the answer string itself, is structurally unable to make. The limitations
below give the trace record that demonstrates it.

One finding did not require the benchmark, because it surfaced while preparing
to run it. It appears in the next section, being as much limitation as result.

\section{Limitations}

\textbf{Sample size, and an incomplete grid.} Results come from 250-question
stratified samples of each dev split rather than the full splits. At that size
the per-system 95\% interval on token F1 spans roughly twelve points ---
$[0.451, 0.571]$ for \texttt{agentic\_full} on HotpotQA --- and an earlier
draft argued from the overlap of such intervals that the deficit against
\texttt{self\_ask} could not be separated from noise. That was the wrong test:
the systems answer the same questions, the comparison is paired, and the
paired interval on the difference, $[-0.145, -0.039]$ on HotpotQA, excludes
zero. The width of the marginal intervals is still the right caution about
effect sizes --- nothing here should be read as a precise one --- but
significance is decided by the paired column and not by overlap. The grid is
also not finished. The ablations run against the same 250-question evaluation
split rather than a subsample; at the time of writing the six baseline and
full-system configurations are complete on both datasets and all three
ablations are complete on HotpotQA, while on 2WikiMultihopQA
\texttt{agentic\_no\_verifier} was still in flight when the tables were last
generated and \texttt{agentic\_no\_planner} and \texttt{agentic\_no\_kg} had
not started. Every ablation conclusion in this report is therefore a
one-dataset conclusion: the backward edge's contribution is a single positive
number, the planner's a single null pointing the wrong way, and the graph's a
single null. Each is marked so where it is due, and none of the three is
general until the 2WikiMultihopQA column exists.

\textbf{Model scale.} Every result is conditional on a 4-bit quantised 8B model.
The planner findings in particular are unlikely to transfer: a frontier model
would probably label its own strategy correctly, making the DAG-derived override
unnecessary, and would produce better decompositions in the first place,
reducing the headroom the feedback loop has to work in. The cost profile is
equally hardware-specific; at 52 tokens per second the case for replacing model
calls with rules is far stronger than against a fast hosted endpoint.

\textbf{Domain.} Both benchmarks are built from Wikipedia: clean, well-segmented,
densely hyperlinked, and written in a consistent register --- and the entity
graph exploits exactly those properties. Nothing in this work is evidence that
it generalises to legal, biomedical, or enterprise collections, where entity
mentions are less standardised and passages are longer and less
self-contained.

\textbf{Entity linking.} Entities are linked by alias matching against a table
built from passage titles and their mechanical variants, not by a trained
linker. This is fast, deterministic, and interpretable, and it fails silently on
nominal and pronominal reference and on any entity whose surface form in running
text does not resemble its page title. Falling back to retrieved passage titles
as seeds limits the damage but does not measure it.

\textbf{An uncalibrated threshold.} The Verifier's confidence threshold of 0.55
and the weights of the confidence blend were set a priori, not learned. The
threshold directly controls the re-plan rate, itself a headline agent metric, so
tuning it on the evaluation questions would be leakage. The disjoint
50-question calibration slice of each dataset exists for exactly this purpose:
the design called for the threshold to be swept over it and frozen before the
evaluation run rather than accepted as a guess. That is not what happened.
The sweep reported in Tables~\ref{tab:calibration}
and~\ref{tab:calibration-reliability} was run after the evaluation, over the
range $[0.40, 0.75]$ in steps of 0.01, on the 50-question HotpotQA calibration
slice run \texttt{agentic\_full\_hotpotqa\_20260904T184256Z}, reading
confidences that a single run at 0.55 had already produced; the threshold
that governed every evaluation run is the a priori one. The sweep returned a
more useful answer than a better threshold. The confidence blend is not a
probability. Expected calibration error is 0.332 and maximum calibration error
0.636 against a base rate of 0.340, and the Brier score of 0.339 is
\emph{worse} than the constant predictor that always returns the base rate, a
Brier skill of $-0.510$.
Nothing in this report should be read as saying that a confidence of 0.7 means
a seven-in-ten chance of being correct, and the report does not say it
anywhere. What the same signal does have is weak discrimination: the area under
the ROC curve is 0.633, above chance but only modestly. Calibration and
discrimination are different properties and this signal has the second without
the first, which is exactly the profile of a usable \emph{ranking cut-point} and
an unusable probability. The threshold should be understood, and is used, as
the former.

On the threshold itself the sweep declines to act, and the reasons matter more
than the value. Youden's $J$ is maximised at 0.54, but a bootstrap over the
selection returns a 95\% interval of $[0.45, 0.64]$ --- a width of 0.19 over a
range swept from 0.40 to 0.75, which does not identify an operating point at
all. The
configured 0.55 lies inside that interval, and the paired comparison of
decision accuracy between 0.54 and 0.55 gives $[-0.060, +0.060]$ at $p = 1.000$
with the selected value enjoying the advantage of having been fitted on the
same 50 questions it was then scored on. Fifty labelled questions is below the
threshold at which the calibration module will propose replacing a configured
value at all, so its recommendation to retain 0.55 is a refusal to act on thin
evidence rather than an endorsement of the guess. The honest statement is that
the threshold remains unvalidated and that the study now knows precisely how
unvalidated it is. A genuine sweep would require one full evaluation run per
candidate threshold, since changing the threshold changes whether a re-plan
fires and therefore changes the plan, the evidence and the answer; the
post-hoc sweep reads confidences that were all produced under 0.55 and cannot
substitute for that.

This connects directly to the Verifier's 92 and 154 false accepts. With
discrimination this weak, moving the acceptance boundary trades false accepts
for false rejects along a shallow curve rather than removing them, so the
permissiveness reported in the previous section is a property of the confidence
signal and not of where the line was drawn on it. Improving the loop means
improving the signal --- atomic claim decomposition, cross-source agreement, a
better entailment model --- not retuning the threshold.

\textbf{The Verifier checks the sentence, not the answer.} The NLI hypothesis
is the Synthesizer's \texttt{answer\_sentence} field, chosen because a
declarative sentence is what an entailment model is trained on; the
\texttt{answer} string that the metrics score is never entailed against
anything. The two can come apart, and when they do the blend cannot see it.
HotpotQA question \texttt{5a7af32e55429931da12c99c} is the cleanest instance
in the trace. The system answered \emph{New York City} where the gold answer
is \emph{Brooklyn, New York}, and emitted as its answer sentence a sentence
copied from an evidence passage unrelated to the question. The Verifier
scored NLI support at 0.995 and citation grounding at 1.0, since the copied
sentence is trivially entailed by the sentence it was copied from, recorded a
retrieval agreement of 0.0, and accepted the answer at a confidence of 0.82.
Every component signal was computed correctly; the architecture asked the
wrong question. A verifier that accepts an answer because its
\emph{explanation} is well supported will accept any answer whose explanation
is a quotation, and that is a structural gap rather than a tuning one. The
correction --- entail a hypothesis built from the question and the answer
string, or check that the answer string occurs in the entailed sentence ---
is straightforward, and it was not in place for any run reported here.

\textbf{Re-planning overwrites the previous cycle's state.} Sub-query
identifiers restart at \texttt{q1} in every plan revision, and the per-sub-query
result, knowledge-graph result and extracted-answer maps on the question state
are keyed by that identifier. A re-plan's \texttt{q1} therefore overwrites
cycle~0's \texttt{q1}, while any identifier the new plan does not reuse
survives from the earlier cycle, stale. Of the 118 HotpotQA questions on which
\texttt{agentic\_full} re-planned, 115 lost their cycle-0 retrieval results in
this way and 92 carry at least one stale identifier from a plan that was no
longer current. Two consequences follow. The retrieval metrics reported for
the 47.2\% of HotpotQA questions that re-planned are computed over a
last-cycle-plus-stale pool rather than over everything the question retrieved,
so the agentic rows of the retrieval tables are a lower bound on what was
found. And the trace is not the complete audit record the design document
promises: the plans are preserved immutably, but the per-sub-query results
they produced are not, so a reader cannot reconstruct what the Synthesizer
was shown in cycle~0 of a re-planned question. The fix is to namespace the
state by cycle; it is recorded as a known defect in the design document and
was not applied to any run reported here.

\textbf{The wall-clock budget is cooperative.} The 300\,s question budget is
tested between DAG nodes and before a model call is issued, never during one,
and the NLI and reranker invocations carry no timeout at all. It therefore
cannot pre-empt a call that blocks, which is how one HotpotQA question ran
for 18.3 hours inside a single Synthesizer call --- almost certainly a
suspended host rather than a hung model --- without any guard firing. The
per-question latency tables print the median with that maximum beside it so
the outlier stays visible; the defect stands, and a per-invocation bound is
the correction.

\textbf{Silent degradation, discovered the hard way.} The Verifier never blocks:
if the NLI cross-encoder cannot be loaded it falls back to token containment
between the answer and the cited evidence. That fallback is safe for the process
and catastrophic for the results. On a comparison claim supported by only one
operand's date, token containment scored 1.000 support where DeBERTa scores
0.0008 entailment on the same pair. Because NLI carries 0.45 of the confidence
blend, a silent fallback pushes essentially every answer above the 0.55
threshold, every verdict becomes \texttt{accept}, and the re-plan loop --- the
mechanism the project exists to measure --- never fires. The run completes, the
numbers look plausible, and the central result is an artefact of a missing model
file. A preflight check now confirms the encoder scores a self-entailment
correctly before any evaluation starts. The general lesson outweighs the bug:
graceful degradation and honest measurement are in direct tension, because the
fallbacks that keep a long run alive are exactly the ones that quietly change
what is measured. Every degradation path needs a preflight gate or a trace flag.

\section{Future Work}

\textbf{Multi-agent coordination beyond a fixed topology.} The agent graph here
is static: the orchestrator invokes the same four agents in the same order, and
the only dynamism is the re-plan edge. A natural extension instantiates agents
per question --- several retrieval agents specialising in different sub-corpora,
or a negotiation between graph-first and text-first strategies whose
disagreement is itself a signal. Workflow-level agent composition has a
precedent in the geospatial copilot literature~\cite{lee2025multiagent}, and how
many agents such a scheme can add before cost outruns benefit is directly
measurable with the instrumentation already in place.

\textbf{Learned rather than heuristic decision-making.} The tool-routing table,
the rerank gate, and the confidence blend are hand-specified functions whose
inputs are logged on every question, making each a supervised learning problem
with a free training set. The interesting version is not replacing the rules
with a classifier but quantifying the gap: if a learned router beats the rule
table by a negligible margin, that is a stronger result for the rules than any
argument. The first step is cheaper than learning anything: the trace records
the Planner's hint and the lexical features on every sub-query, so the
agreement between the hint that was honoured and the rule that would have
fired can be computed offline, and it would say whether six of the seven rules
were redundant or merely unreached.

\textbf{Stronger fact checking.} The Verifier checks whether the final answer
sentence is entailed by its citations. It does not check intermediate hops, does
not cross-reference independent sources, and cannot distinguish an unsupported
claim from one whose support was simply not retrieved. A verifier that
decomposed the answer into atomic claims, verified each independently, and
located contradictions between sources would give the re-plan directive far more
precise information about what to look for next --- and directive quality, not
the loop's existence, is plausibly what limits its usefulness.

\textbf{Domain transfer.} The brief suggests legal and biomedical retrieval, and
each stresses a different part. Biomedical corpora have curated ontologies, so
the KG Navigator could traverse an authored graph instead of an induced one,
isolating how much of the graph's value comes from structure rather than from
the induction procedure. Legal corpora invert the challenge: citation networks
are explicit and reliable, but passages are long, densely cross-referential, and
full of terms of art whose meaning is document-specific, testing the
alias-matching linker to destruction.

\section{Toward a Thesis Extension}
\label{sec:thesis}

Two of the thesis directions in the assignment brief follow from this system
rather than requiring a new one. Both would also be the natural place to
introduce the real-time external tools the brief describes and this system
deliberately omits: the registry accepts a connector as one more entry, and
the budgeting discipline applied here to model calls is what such a connector
would need.

The self-reflective fact-checking track is the closer fit. The Verifier already
scores confidence and triggers further retrieval when it is low, which is the
mechanism the brief describes; what it lacks is multi-source cross-verification,
provenance annotation, and calibrated uncertainty reported to the user rather
than consumed internally. Extending it to a multimodal setting --- validating
retrieved imagery against textual reports, as the brief's flood-extent example
illustrates --- needs a cross-modal entailment model in place of the NLI stage,
but leaves the orchestration, the directive schema, and the anti-loop machinery
intact.

The agentic query-decomposition track is the other. This Planner decomposes a
static text corpus into a graph of at most five nodes; the spatio-temporal
version decomposes over archives indexed by time and location, where sub-queries
carry temporal extents and dependencies encode change detection between slices.
The structure --- validated DAG, placeholder substitution, per-node tool
routing, verification with feedback --- is unchanged; the tool registry and the
cost model are not. The cost model is where the work sits: when one tool
invocation is a satellite-imagery request rather than a BM25 query, the
budgeting discipline applied here to language-model calls becomes the central
design problem rather than an implementation detail.
